% paper3_intercropping.tex -- Paper 3. Typeset alone with this file, or
% inside the volume through main.tex.
\documentclass[../main.tex]{subfiles}
\begin{document}
\begin{paper}{Maize--Cowpea Intercropping and the Land Equivalent Ratio}{Maize--cowpea intercropping}
\label{paper:intercropping}

\begin{paperabstract}
Growing a cereal and a legume together on the same land can produce more in
total than growing each crop alone on the same area. This paper compares
maize and cowpea grown as sole crops with the two crops grown together in
alternating rows, at three sites. The advantage of the intercrop is measured
by the land equivalent ratio, calculated from the intercrop yield of each
species relative to its sole-crop yield, and the contribution of each crop to
the combined advantage is reported through its partial ratio. All yields are
synthetic; the paper demonstrates how the land equivalent ratio is calculated,
tabulated and interpreted.
\end{paperabstract}
\paperkeywords{intercropping; maize; cowpea; land equivalent ratio; cereal--legume systems}

\section{Introduction}

Intercropping, the growing of two or more crops together for at least part
of their life cycle, is an old practice that remains widespread in low-input
farming systems. It has attracted renewed interest because it can raise total
production without additional inputs, or make production more stable, and so
offers one route towards the sustainable intensification of agriculture
\parencite{brooker2015}. The advantage arises when the component crops use
resources in complementary ways, for example by rooting at different depths,
by drawing on light at different times, or, in the case of a legume, by fixing
atmospheric nitrogen that the cereal would otherwise have to take from the
soil.

Comparing an intercrop with sole crops is not straightforward, because the
intercrop produces two harvests while each sole crop produces one. The land
equivalent ratio resolves this by asking how much land would be needed under
sole cropping to produce the yields obtained from one unit of land under
intercropping. It is the sum, over the component crops, of the intercrop yield
of each crop divided by its sole-crop yield \parencite{mead1980}. A ratio
greater than one shows that the intercrop uses land more efficiently than the
sole crops; a ratio of 1.2, for example, means that 20 per cent more land
would be needed to match the intercrop by growing the crops separately.

The value of the ratio depends on the sole-crop yields used to standardise
it, and the choice of those yields must be stated when results are reported
\parencite{mead1980}. In this trial each sole crop is grown in the same blocks
as the intercrop, so that the standardising yields come from the same site,
season and management. Cereal--legume combinations are also a component of
integrated soil fertility management, in which legumes contribute to the
nitrogen economy of the following cereal \parencite{vanlauwe2010}.

The aims of this paper are to report the sole-crop and intercrop yields of
maize and cowpea at three sites, to calculate the partial and total land
equivalent ratios, and to identify the sites at which the intercrop gives a
yield advantage.

\input{\subfix{methods/paper3_methods}}

\input{\subfix{results/paper3_results}}

\end{paper}
\end{document}
